\chapter{Continuidad, densidad y uniformidad}\label{cap:nuevo02}\label{chap:05}\label{chap:06}\label{chap:07}

La topología métrica permite precisar qué significa que una función respete la proximidad. En este capítulo estudiaremos primero la continuidad local y sus diferentes criterios; después veremos cómo los conjuntos densos y las bases numerables organizan las aproximaciones; finalmente distinguiremos el control local del control uniforme, que preserva sucesiones de Cauchy y permite extender funciones a puntos límite. Se integran las lecturas 5--7 de \cite{valencia}. Usaremos las propiedades de distancias, clausuras y subespacios del capítulo~\ref{cap:nuevo01}; los resultados posteriores se apoyarán en esos criterios y en el control explícito de los cuantificadores.

\section{Controlar la imagen de puntos próximos}\label{lectura:05}
Para estudiar la continuidad hay que distinguir la \emph{imagen} $f(A)$ de la \emph{imagen inversa} $f^{-1}(B)$. La primera se aplica a subconjuntos del dominio y la segunda a subconjuntos del codominio; la notación de imagen inversa no exige que $f$ sea invertible. Con esta distinción, conectaremos las formulaciones mediante bolas y tolerancias, abiertos y sucesiones.

\begin{definicion}\label{ch05:continua}\index{continuidad!en un punto}
Sean $(X,d)$ y $(Y,\rho)$ espacios métricos. Una función $f:X\to Y$ es continua en $x\in X$ si para toda vecindad $V$ de $f(x)$, el conjunto $f^{-1}(V)$ es vecindad de $x$.
\end{definicion}

\begin{teorema}[Criterio de tolerancias]\label{ch05:epsdelta}
La función $f$ es continua en $x$ si y solo si
\[
\forall\eps>0\ \exists\delta>0\ \forall u\in X:
\quad d(x,u)<\delta\ \Longrightarrow\ \rho(f(x),f(u))<\eps.
\]
\end{teorema}
\begin{proof}
Si $f$ es continua en $x$, la imagen inversa de la vecindad $B_Y(f(x),\eps)$ contiene alguna bola $B_X(x,\delta)$. La inclusión de bolas es exactamente la implicación del enunciado. Recíprocamente, sea $V$ una vecindad de $f(x)$ y elijamos $\eps>0$ con $B_Y(f(x),\eps)\subset V$. La condición proporciona $\delta>0$ con $B_X(x,\delta)\subset f^{-1}(B_Y(f(x),\eps))\subset f^{-1}(V)$, así que esta última es vecindad de $x$.
\end{proof}
El orden de los cuantificadores explica el significado: primero se prescribe una precisión en las imágenes; después se elige una tolerancia en el dominio que sirve para \emph{todos} los puntos suficientemente próximos. La tolerancia puede depender del punto $x$; exigir una sola tolerancia para todos los puntos conduce más adelante a continuidad uniforme.

\begin{teorema}[Criterios globales]\label{ch05:criterios}\index{continuidad!criterios topologicos@criterios topológicos}
Para $f:X\to Y$, son equivalentes:
\begin{enumerate}[label=(\roman*)]
\item $f^{-1}(U)$ es abierto en $X$ para todo abierto $U$ en $Y$.
\item Para todo $A\subset X$, $f(\cl A)\subset\cl{f(A)}$.
\item $f^{-1}(F)$ es cerrado en $X$ para todo cerrado $F$ en $Y$.
\item $f$ es continua en cada punto de $X$.
\end{enumerate}
Cuando se cumplen, se dice que $f$ es continua en $X$.
\end{teorema}
\begin{proof}
\emph{(i)$\Rightarrow$(ii).} Si $x\in\cl A$ y $U$ es abierto con $f(x)\in U$, entonces $f^{-1}(U)$ es un abierto que contiene a $x$. Por ello intersecta a $A$, y la imagen de un punto de esa intersección pertenece a $U\cap f(A)$. Como toda bola alrededor de $f(x)$ es uno de esos abiertos, $f(x)\in\cl{f(A)}$.

\emph{(ii)$\Rightarrow$(iii).} Sea $F$ cerrado y $E=f^{-1}(F)$. La inclusión $f(E)\subset F$ implica $\cl{f(E)}\subset F$. Si $x\in\cl E$, por (ii), $f(x)\in f(\cl E)\subset\cl{f(E)}\subset F$, de modo que $x\in E$. Así $E$ es cerrado.

\emph{(iii)$\Rightarrow$(i).} Para $U$ abierto, $Y\setminus U$ es cerrado y
\[
X\setminus f^{-1}(U)=f^{-1}(Y\setminus U)
\]
es cerrado. Por tanto $f^{-1}(U)$ es abierto.

\emph{(i)$\Rightarrow$(iv).} Si $V$ es vecindad de $f(x)$, contiene una bola abierta $U$ centrada en $f(x)$. Por (i), $f^{-1}(U)$ es abierto y contiene a $x$; contiene entonces una bola alrededor de $x$. Como $f^{-1}(U)\subset f^{-1}(V)$, esta última es vecindad de $x$.

\emph{(iv)$\Rightarrow$(i).} Sea $x\in f^{-1}(U)$, con $U$ abierto. Este $U$ es vecindad de $f(x)$, así que $f^{-1}(U)$ es vecindad de $x$. Lo es para cada uno de sus puntos y, por tanto, es abierto.
\end{proof}
No se afirma en (ii) igualdad. Por ejemplo, $f(t)=\arctan t$ es continua en $\R$, pero para $A=\R$ su imagen es $(-\pi/2,\pi/2)$ y la clausura de esa imagen contiene además los dos extremos.

\section{Ejemplos y reglas de construcción}
Toda función cuyo dominio lleva la topología discreta es continua, porque cualquier imagen inversa es un subconjunto abierto del dominio. La identidad de un espacio en sí mismo es continua, ya que sus imágenes inversas no cambian los conjuntos.

Si $f:X\to Y$ y $g:Y\to Z$ son continuas, también lo es $g\circ f$: para un abierto $U\subset Z$,
\[
(g\circ f)^{-1}(U)=f^{-1}(g^{-1}(U)),
\]
y primero la continuidad de $g$ y luego la de $f$ garantizan apertura. Se está usando la preimagen de conjuntos, no una inversa funcional de $f$ o $g$.

Una función $L$-Lipschitz, es decir, una función con $\rho(f(x),f(y))\le Ld(x,y)$, es continua. Si $L>0$, basta tomar $\delta=\eps/L$; si $L=0$, todas las imágenes coinciden y cualquier $\delta>0$ sirve. Así, la función $x\mapsto d(x,A)$ para $A\ne\varnothing$ es continua por el ejercicio~\ref{ej:lectura-3-7}, y la norma $v\mapsto\norm v$ es continua porque
\[
|\norm u-\norm v|\le\norm{u-v}.
\]
Esta última desigualdad se prueba aplicando la triangular a $u=(u-v)+v$ y luego intercambiando $u,v$.

La inclusión $Y\hookrightarrow X$ de un subespacio es continua: la preimagen de un abierto ambiental es su intersección con $Y$, abierta relativamente. Una restricción $f|_Y$ de una función continua es continua por composición con esa inclusión. Además, si una función toma valores en un subespacio $Z\subset Y$, es continua como función a $Z$ si y solo si lo es como función a $Y$. En una dirección se compone con la inclusión; en la otra, para $U=Z\cap G$ abierto en $Z$, se tiene $f^{-1}(U)=f^{-1}(G)$ porque toda imagen pertenece a $Z$. Estos hechos incorporan el resultado sobre subespacios que aparece repetido en la lectura original.

\begin{teorema}[Criterio secuencial]\label{ch05:secuencial}\index{continuidad!criterio secuencial}
Una función entre espacios métricos es continua en $x$ si y solo si $x_n\to x$ implica $f(x_n)\to f(x)$.
\end{teorema}
\begin{proof}
Supongamos continuidad y fijemos $\eps>0$. Elegimos $\delta>0$ según el teorema~\ref{ch05:epsdelta}. Como $x_n\to x$, hay $N$ tal que $d(x_n,x)<\delta$ para $n\ge N$. Entonces $\rho(f(x_n),f(x))<\eps$ para esos índices.

Para el recíproco usamos la negación exacta de continuidad: existe $\eps_0>0$ tal que para \emph{todo} $\delta>0$ existe $u\in X$ con $d(u,x)<\delta$ y $\rho(f(u),f(x))\ge\eps_0$. Tomando $\delta=1/n$, elegimos $x_n$ con esas propiedades. Entonces $x_n\to x$, pero ninguna imagen entra en la bola $B(f(x),\eps_0)$. Esto contradice el criterio secuencial. La desigualdad de la negación es $\ge\eps_0$, y no es necesario reemplazarla por una estricta.
\end{proof}

\section{Funciones abiertas, cerradas y homeomorfismos}
\begin{definicion}\index{funcion@función!abierta}\index{funcion@función!cerrada}\index{homeomorfismo}
Una función $f:X\to Y$ es abierta si $f(U)$ es abierto para cada abierto $U\subset X$, y cerrada si $f(F)$ es cerrado para cada cerrado $F\subset X$. Es un homeomorfismo si es biyectiva, continua y su inversa $f^{-1}:Y\to X$ es continua. En este caso se escribe $X\cong Y$.
\end{definicion}
Estas condiciones sobre imágenes directas son diferentes de continuidad. La identidad de $\R$ con su métrica discreta hacia $\R$ usual es continua, pero no es abierta, pues la imagen de $\{0\}$ no es abierta, ni cerrada, pues la imagen del conjunto discreto cerrado $\Q$ no es cerrada en la recta usual. La identidad en la dirección contraria es abierta y cerrada, porque todo subconjunto del codominio discreto es abierto y cerrado; sin embargo, no es continua.

\begin{ejemplo}\label{ch05:homeorecta}
La función $h:\R\to(-1,1)$ dada por $h(t)=t/(1+|t|)$ es un homeomorfismo. Su valor tiene módulo menor que $1$. Para $s\in(-1,1)$, definamos $k(s)=s/(1-|s|)$. Como $|k(s)|=|s|/(1-|s|)$,
\[
h(k(s))=\frac{s/(1-|s|)}{1+|s|/(1-|s|)}=s.
\]
Del mismo modo, $1-|h(t)|=1/(1+|t|)$ implica $k(h(t))=t$. Por tanto son funciones inversas y ambas son continuas por las reglas elementales de cálculo: los denominadores son positivos en sus respectivos dominios. Esto demuestra biyectividad y continuidad de las dos direcciones, no solamente una semejanza geométrica.
\end{ejemplo}

\section{Ejercicios sobre continuidad y homeomorfismos}
\begin{ejercicio}[Lectura 5, ejercicio 1]\label{ej:lectura-5-1}
Demostrar que cada una de las siguientes condiciones equivale a continuidad de $f:X\to Y$:
\begin{enumerate}[label=(\roman*)]
\item Para todo $A\subset Y$, $f^{-1}(\operatorname{int}A)\subset\operatorname{int}(f^{-1}(A))$.
\item Para todo $A\subset Y$, $\cl{f^{-1}(A)}\subset f^{-1}(\cl A)$.
\end{enumerate}
En (ii), se corrige la letra $B$, no definida en el original, por $A$.
\end{ejercicio}
\begin{solucion}
\begin{enumerate}[label=(\roman*)]
\item Si $f$ es continua, $f^{-1}(\operatorname{int}A)$ es abierto y está contenido en $f^{-1}(A)$. El interior de un conjunto contiene a todos sus subconjuntos abiertos; de ahí la inclusión. Si la condición vale y $U$ es abierto, entonces $\operatorname{int}U=U$ y $f^{-1}(U)\subset\operatorname{int}(f^{-1}(U))$. La inclusión inversa es automática, de modo que $f^{-1}(U)$ es abierto y $f$ continua.
\item Si $f$ es continua, $f^{-1}(\cl A)$ es cerrado y contiene a $f^{-1}(A)$; por ello contiene a su clausura. Recíprocamente, para un cerrado $F$ tenemos $\cl F=F$, y la condición da $\cl{f^{-1}(F)}\subset f^{-1}(F)$. Por tanto todas las preimágenes de cerrados son cerradas y $f$ es continua.
\end{enumerate}
\end{solucion}

\begin{ejercicio}[Lectura 5, ejercicio 2]\label{ej:lectura-5-2}
Demostrar que una función es continua en el espacio si y solo si es continua en cada punto.
\end{ejercicio}
\begin{solucion}
Las dos direcciones están desarrolladas en el teorema~\ref{ch05:criterios}. El mecanismo es local: la preimagen de un abierto $U$ es abierta exactamente cuando contiene una bola alrededor de cada $x$ con $f(x)\in U$. La continuidad en ese $x$ produce esa bola, y la continuidad global garantiza que la preimagen de una bola alrededor de $f(x)$ es un abierto que contiene a $x$. Así se recupera en cada dirección la definición requerida.
\end{solucion}

\begin{ejercicio}[Lectura 5, ejercicio 3]\label{ej:lectura-5-3}
Probar que $f$ es continua si y solo si las preimágenes de todas las bolas abiertas son abiertas. Dar una función no continua cuyas preimágenes de todas las bolas cerradas sean cerradas.
\end{ejercicio}
\begin{solucion}
La necesidad sigue de que las bolas abiertas son abiertas. Para la suficiencia, todo abierto $U$ del codominio es una unión $\bigcup_iB_i$ de bolas abiertas. Como las preimágenes preservan uniones,
\[
f^{-1}(U)=\bigcup_i f^{-1}(B_i),
\]
que es abierto por hipótesis.

Para la segunda parte, tomemos $X=\R$ con la métrica usual, $Y=\R$ con la discreta y $f(t)=t$. Una bola cerrada de radio $r<1$ en $Y$ es un singleton, y una de radio $r\ge1$ es todo $Y$. Sus preimágenes son singletons o $\R$, todos cerrados en $X$. Pero $\{0\}$ es abierto en $Y$ y su preimagen no es abierta en $X$; así $f$ no es continua. El argumento con bolas abiertas no se transfiere a bolas cerradas, porque una unión arbitraria de cerrados no tiene por qué ser cerrada.
\end{solucion}

\begin{ejercicio}[Lectura 5, ejercicio 4]\label{ej:lectura-5-4}
Sean $f,g:X\to\R$ continuas. Probar que (i) $\{x:f(x)<g(x)\}$ es abierto y (ii) $\{x:f(x)\le g(x)\}$ es cerrado. Se corrige el dominio de $g$ a $X$, necesario para que ambas expresiones estén definidas.
\end{ejercicio}
\begin{solucion}
\begin{enumerate}[label=(\roman*)]
\item Si $f(x)<g(x)$, sea $\eta=(g(x)-f(x))/3>0$. La continuidad de ambas funciones permite elegir $\delta>0$ tal que $d(x,y)<\delta$ implique $|f(y)-f(x)|<\eta$ y $|g(y)-g(x)|<\eta$. Entonces
\[
f(y)<f(x)+\eta<g(x)-\eta<g(y).
\]
La bola está contenida en el conjunto, que es abierto.
\item El complemento es $\{x:g(x)<f(x)\}$, abierto por (i) aplicado con los papeles intercambiados. Por tanto el conjunto original es cerrado.
\end{enumerate}
\end{solucion}

\begin{ejercicio}[Lectura 5, ejercicio 5]\label{ej:lectura-5-5}\index{extension@extensión!por cero}\index{frontera}
Sea $A\subset X$ cerrado y $f:A\to\R$ continua, con $f=0$ en la frontera $\partial A=\cl A\cap\cl{X\setminus A}$. Demostrar que la extensión $g:X\to\R$, igual a $f$ en $A$ y a cero fuera de $A$, es continua.
\end{ejercicio}
\begin{solucion}
Distinguimos tres posiciones del punto $x$.

Si $x\notin A$, como $A$ es cerrado existe una bola alrededor de $x$ contenida en $X\setminus A$. Allí $g$ es constante y es continua en $x$.

Si $x\in\operatorname{int}A$, existe $r>0$ con $B_X(x,r)\subset A$. Dado $\eps>0$, la continuidad relativa de $f$ da $\delta_0>0$ tal que $y\in A$ y $d(x,y)<\delta_0$ implican $|f(y)-f(x)|<\eps$. Tomando $\delta=\min\{r,\delta_0\}$ se obtiene la misma desigualdad para $g$ y todo $y\in X$ con $d(x,y)<\delta$.

Finalmente, si $x\in A\setminus\operatorname{int}A$, todo entorno de $x$ intersecta al complemento de $A$, así que $x\in\partial A$ y $g(x)=f(x)=0$. Por continuidad relativa existe $\delta>0$ con $|f(y)|<\eps$ para $y\in A\cap B(x,\delta)$. Para $y$ de esa bola fuera de $A$, $g(y)=0$; para los puntos dentro de $A$, $|g(y)|=|f(y)|<\eps$. La misma tolerancia sirve para ambos lados de la frontera. Por tanto $g$ es continua en todos los puntos.
\end{solucion}

\begin{ejercicio}[Lectura 5, ejercicio 6: corrección del recíproco]\label{ej:lectura-5-6}\index{soporte}
Sea $S=\supp(f)=\cl{\{x:f(x)\ne0\}}$. Examinar la afirmación «$f$ es continua si y solo si $f=0$ en $\partial S$» y establecer una versión válida.
\end{ejercicio}
\begin{solucion}
La implicación directa es verdadera. Si $x\in\partial S$ y $f(x)\ne0$, por continuidad hay una bola en la que $|f(y)-f(x)|<|f(x)|/2$, así que $|f(y)|>|f(x)|/2>0$. Esa bola está contenida en $S$, lo que convierte a $x$ en punto interior de $S$ y contradice $x\in\partial S$.

El recíproco es falso. En $\R$, tomemos $f(x)=1$ para $x<0$ y $f(x)=2$ para $x\ge0$. La función no es continua en $0$, pero nunca se anula, de modo que $S=\R$ y $\partial S=\varnothing$. La condición en la frontera se satisface de forma vacía y no controla el salto interior.

Una versión exacta es: $f$ es continua en $X$ si y solo si $f|_S$ es continua en el subespacio $S$ y $f=0$ en $\partial S$. La necesidad se sigue de la continuidad de las restricciones y de la prueba anterior. Para la suficiencia, $S$ es cerrado por definición, $f$ es cero en $X\setminus S$ y el ejercicio~\ref{ej:lectura-5-5} aplica a la función continua $f|_S$. Su extensión por cero es precisamente $f$.
\end{solucion}

\begin{ejercicio}[Lectura 5, ejercicio 7]\label{ej:lectura-5-7}
Probar: (i) $f$ es abierta si y solo si las imágenes de las bolas abiertas son abiertas; (ii) si $f$ es biyectiva, es abierta si y solo si es cerrada; (iii) construir una función suprayectiva abierta que no sea cerrada y una cerrada que no sea abierta.
\end{ejercicio}
\begin{solucion}
\begin{enumerate}[label=(\roman*)]
\item La necesidad es inmediata porque una bola abierta es un abierto. Para la suficiencia, escribimos un abierto $U$ como unión de bolas $B_i$; entonces $f(U)=\bigcup_if(B_i)$ es una unión de abiertos.
\item Para una biyección se cumple $f(X\setminus A)=Y\setminus f(A)$. Si $f$ es abierta y $F$ cerrado, el conjunto $Y\setminus f(F)=f(X\setminus F)$ es abierto, así que $f(F)$ es cerrado. Si $f$ es cerrada y $U$ abierto, el conjunto $Y\setminus f(U)=f(X\setminus U)$ es cerrado, de modo que $f(U)$ es abierto. La identidad de complementos utiliza tanto inyectividad como suprayectividad.
\item La proyección $\pi:\R^2\to\R$, $\pi(x,y)=x$, es suprayectiva y abierta. Si $a\in\pi(U)$ con $U$ abierto, existe $b$ con $(a,b)\in U$ y un disco de radio $r>0$ centrado en $(a,b)$ contenido en $U$. Entonces $(a-r,a+r)\subset\pi(U)$, pues $(t,b)$ pertenece a ese disco cuando $|t-a|<r$. Sin embargo, el conjunto $F=\{(x,y):xy=1\}$ es cerrado, como preimagen de $\{1\}$ por la función continua producto, y $\pi(F)=\R\setminus\{0\}$ no es cerrado.

Para una función incluso suprayectiva y cerrada que no sea abierta, tomemos $h:\R\to[0,\infty)$, $h(t)=\max\{t,0\}$. Si $F$ es cerrado, su imagen es $F\cap[0,\infty)$, con el punto $0$ añadido cuando $F$ intersecta $(-\infty,0)$. En ambos casos es un cerrado de $[0,\infty)$. Pero $h((-2,-1))=\{0\}$ no es abierto en $[0,\infty)$: toda bola relativa alrededor de $0$ contiene puntos positivos.
\end{enumerate}
\end{solucion}

\begin{ejercicio}[Lectura 5, ejercicio 8]\label{ej:lectura-5-8}\index{pegado!sobre abiertos}
Sea $(A_i)_{i\in I}$ una cubierta abierta de $X$ y sean $f_i:A_i\to Y$ continuas que coinciden en todas las intersecciones. Probar que la función obtenida pegándolas es continua.
\end{ejercicio}
\begin{solucion}
Para $x\in X$, elijamos $i$ con $x\in A_i$ y definamos $g(x)=f_i(x)$. Si también $x\in A_j$, la hipótesis da $f_i(x)=f_j(x)$; por tanto la definición no depende de la elección. Cualquier función con las restricciones prescritas tiene necesariamente este valor en cada $x$, de modo que es única.

Sea $U$ abierto en $Y$. Entonces
\[
g^{-1}(U)=\bigcup_{i\in I}f_i^{-1}(U).
\]
Cada preimagen es abierta en $A_i$ por continuidad. Como $A_i$ es abierto en $X$, el ejercicio~\ref{ej:lectura-4-2} garantiza que también es abierta en $X$. Su unión es abierta y el criterio global prueba continuidad de $g$. No se requiere que la cubierta sea finita, porque las uniones arbitrarias de abiertos preservan apertura.
\end{solucion}

\begin{ejercicio}[Lectura 5, ejercicio 9]\label{ej:lectura-5-9}
Demostrar que ser homeomorfos es una relación de equivalencia entre espacios métricos.
\end{ejercicio}
\begin{solucion}
La identidad de $X$ es una biyección continua cuya inversa es ella misma; esto prueba reflexividad. Si $h:X\to Y$ es un homeomorfismo, su inversa $h^{-1}:Y\to X$ es biyectiva y continua y tiene inversa continua $h$; esto prueba simetría. Si $h:X\to Y$ y $k:Y\to Z$ son homeomorfismos, $k\circ h$ es una biyección continua y su inversa es $h^{-1}\circ k^{-1}$, continua por composición. Así se prueba transitividad. Cuando se habla de la clase de todos los espacios, «relación de equivalencia» se entiende en ese sentido de clases; sobre cualquier colección que sea un conjunto se obtiene la relación de equivalencia usual.
\end{solucion}

\section{Aproximar un espacio con pocos puntos}\label{lectura:06}
Los criterios de continuidad se han formulado mediante aproximaciones y entornos. Una sucesión solo utiliza un conjunto numerable de términos; cabe preguntar cuándo todo el espacio puede describirse, en términos de aproximación, a partir de un conjunto numerable. La densidad responderá a la primera parte de la pregunta, y la relación entre separabilidad y bases numerables permitirá controlar después familias de abiertos.

Un conjunto es \emph{a lo sumo numerable} si es finito o puede enumerarse mediante números naturales. Usaremos que una unión numerable de conjuntos numerables es numerable y que un producto finito de conjuntos numerables es numerable. Por ejemplo, los pares de naturales pueden enumerarse recorriendo las diagonales $m+n=2,3,\ldots$; una familia numerable de enumeraciones se reduce así a una lista de pares. Las elecciones necesarias se entienden en el marco usual de la teoría de conjuntos con el axioma de elección.

\begin{teorema}\label{ch06:denso}\index{conjunto!denso}
Para $A\subset X$, con $X$ métrico no vacío, son equivalentes:
\begin{enumerate}[label=(\roman*)]
\item $A\ne\varnothing$ y $d(x,A)=0$ para todo $x\in X$.
\item Todo abierto no vacío de $X$ intersecta a $A$.
\item $\cl A=X$.
\end{enumerate}
En este caso se dice que $A$ es denso en $X$.
\end{teorema}
\begin{proof}
Si (i) se cumple y $U$ es abierto no vacío, elijamos $x\in U$ y $r>0$ con $B(x,r)\subset U$. La igualdad $d(x,A)=0$ produce $a\in A$ con $d(x,a)<r$, así que $a\in U$. Si (ii) se cumple, cada bola $B(x,r)$ es un abierto no vacío e intersecta a $A$, lo que implica $x\in\cl A$ para cada $x$. Finalmente, si $\cl A=X$, entonces $A$ es no vacío porque $X$ lo es, y el ejercicio~\ref{ej:lectura-3-7} da $d(x,A)=0$.
\end{proof}

\begin{ejemplo}[Racionales y truncaciones]
En un espacio discreto, cada singleton es abierto. Por tanto un conjunto denso debe contener cada punto, y el único conjunto denso es el espacio entero.

En $\R$, $\Q$ es denso. Si $a<b$, elegimos $n$ con $1/n<b-a$. Existe un entero $m$ tal que $na<m\le na+1$; por ejemplo, $m=\lfloor na\rfloor+1$. Entonces $a<m/n\le a+1/n<b$. Así cada intervalo abierto no vacío contiene un racional.

En $\ell^1$, el subespacio $c_{00}=\operatorname{span}\{e_k:k\ge1\}$ de sucesiones con soporte finito es denso. Para $a\in\ell^1$, su truncación $a^{[N]}=(a_1,\ldots,a_N,0,\ldots)$ satisface
\[
\norm{a-a^{[N]}}_1=\sum_{k>N}|a_k|\longrightarrow0.
\]
La convergencia de la serie justifica que sus colas tiendan a cero. Sin embargo, $c_{00}\ne\ell^1$: la sucesión $(1/n^2)$ pertenece a $\ell^1$ y tiene infinitas coordenadas no nulas. Una familia puede generar un subespacio \emph{denso} sin generar algebraicamente todo el espacio.\index{c00@$c_{00}$}
\end{ejemplo}
Más adelante se demostrará que los polinomios son densos en $C([0,1])$ con la norma uniforme. Esa afirmación necesita un teorema de aproximación, y no se usará como premisa antes de demostrarlo.

\begin{definicion}\index{separabilidad}
Un espacio métrico es separable si posee un subconjunto denso a lo sumo numerable.
\end{definicion}
Así, $\R$ es separable aunque no sea numerable. Un espacio discreto es separable exactamente cuando es numerable, porque su único subconjunto denso es él mismo. No debe confundirse la cantidad de puntos del espacio con la cantidad de puntos necesaria para aproximarlos.

\section{Bases numerables}
\begin{definicion}\index{base!de una topologia@de una topología}\index{segunda numerabilidad}
Una base de la topología de $X$ es una familia $\mathcal B$ de abiertos tal que todo abierto de $X$ es unión de elementos de $\mathcal B$. El espacio es segundo numerable si posee una base a lo sumo numerable.
\end{definicion}
Una formulación equivalente, útil para las pruebas, es: para cada abierto $U$ y cada $x\in U$, existe $B\in\mathcal B$ con $x\in B\subset U$. Si todo abierto es unión de elementos de la base, uno de los que componen $U$ contiene a $x$. En sentido contrario, tomando esos elementos para todos los $x\in U$, su unión es exactamente $U$.

\begin{lema}\label{ch06:baseadenso}
De una base $\mathcal B$ puede obtenerse un conjunto denso cuya cardinalidad no exceda la de $\mathcal B$.
\end{lema}
\begin{proof}
Eliminemos el conjunto vacío de la base, si aparece: no aporta puntos a ninguna unión, por lo que esta sigue siendo base. Para cada $B\in\mathcal B$ elegimos $x_B\in B$, y ponemos $D=\{x_B:B\in\mathcal B\}$. Todo abierto no vacío contiene algún elemento no vacío de la base y, en consecuencia, contiene el punto elegido en él. Por el teorema~\ref{ch06:denso}, $D$ es denso. La aplicación $B\mapsto x_B$ es una sobreyección de la base sobre $D$, lo que da la cota cardinal.
\end{proof}

\begin{teorema}\label{ch06:separablebase}
Un espacio métrico es separable si y solo si es segundo numerable.
\end{teorema}
\begin{proof}
Una base numerable proporciona un conjunto denso numerable por el lema anterior. Para la otra dirección, sea $D=\{a_n:n\in\N\}$ denso; si $D$ es finito no vacío, se pueden repetir sus elementos en una enumeración. Definamos
\[
\mathcal B=\{B(a_n,2^{-k}):n,k\in\N\}.
\]
Esta familia es numerable porque está indexada por pares de naturales. Sea $U$ abierto y $x\in U$; elijamos $\eps>0$ con $B(x,\eps)\subset U$, y luego $k$ con $2^{1-k}<\eps$. La densidad permite elegir $a_n$ con $d(x,a_n)<2^{-k}$. Entonces $x\in B(a_n,2^{-k})$ y, para $y$ en esa bola,
\[
d(x,y)\le d(x,a_n)+d(a_n,y)<2\cdot2^{-k}<\eps.
\]
Por tanto $x\in B(a_n,2^{-k})\subset U$. Esta es la propiedad de base. En el espacio vacío, la familia vacía es una base y el conjunto vacío es denso en él, así que ambas propiedades también coinciden.
\end{proof}
La prueba revela por qué los centros densos son suficientes: se aproxima primero el centro deseado, reservando la mitad de la tolerancia para el radio de la bola. No basta elegir un centro cercano sin controlar también ese radio.

\begin{teorema}[Lindelöf]\label{ch06:lindelof}\index{Lindelof@Lindelöf, teorema de}
Sea $X$ métrico separable y $\mathcal U$ una familia de abiertos. Existe una subfamilia a lo sumo numerable con la misma unión que $\mathcal U$. En particular, toda cubierta abierta de $X$ admite una subcubierta numerable.
\end{teorema}
\begin{proof}
Sea $\mathcal B=\{B_n:n\in\N\}$ una base numerable. Consideremos los índices $n$ para los que $B_n$ está contenido en algún miembro de $\mathcal U$, y para cada uno elijamos un solo $U_n\in\mathcal U$ con $B_n\subset U_n$. La familia elegida es numerable.

Si $x\in\bigcup\mathcal U$, existe $U\in\mathcal U$ con $x\in U$. La propiedad de base proporciona un $B_n$ con $x\in B_n\subset U$. Este índice está entre los seleccionados; aunque el abierto elegido $U_n$ pueda ser distinto de $U$, contiene a $B_n$ y, por tanto, a $x$. Así $\bigcup\mathcal U\subset\bigcup_nU_n$. La inclusión contraria se sigue de que todos los $U_n$ pertenecen a $\mathcal U$. Si la unión es vacía, la subfamilia vacía sirve.
\end{proof}
El teorema no dice que la subcubierta pueda ser finita; esa exigencia define una propiedad más fuerte que estudiaremos bajo el nombre de compacidad.

\section{Ejercicios sobre densidad y numerabilidad}
\begin{ejercicio}[Lectura 6, ejercicio 1]\label{ej:lectura-6-1}
Demostrar que $\ell^\infty$ no es separable.
\end{ejercicio}
\begin{solucion}
Sea $E=\{0,1\}^{\N}\subset\ell^\infty$. Este conjunto no es numerable: si sus elementos se listaran como $a^{(1)},a^{(2)},\ldots$, la sucesión $b$ definida por $b_n=1-a^{(n)}_n$ diferiría de la sucesión número $n$ en la coordenada $n$ y no aparecería en la lista. Es el argumento diagonal de Cantor.

Para $a,b\in E$ distintos, existe una coordenada donde difieren en $1$, y ninguna diferencia tiene módulo mayor que $1$; por tanto $\norm{a-b}_\infty=1$. Las bolas $B(a,1/3)$, $a\in E$, son disjuntas dos a dos: un punto común daría por la triangular una distancia entre centros menor que $2/3$.

Si $D=\{d_n:n\in\N\}$ fuera denso, cada bola contendría algún $d_n$. Asignemos a cada $a\in E$ el primer índice de un punto de $D$ que pertenece a $B(a,1/3)$. Bolas disjuntas reciben índices distintos, lo que produciría una inyección del conjunto no numerable $E$ en $\N$. Es imposible; por tanto $\ell^\infty$ no es separable.
\end{solucion}

\begin{ejercicio}[Lectura 6, ejercicio 2]\label{ej:lectura-6-2}
Demostrar que un espacio normado $V$ es separable si y solo si existe un conjunto a lo sumo numerable $L\subset V$ cuyo subespacio generado sea denso.
\end{ejercicio}
\begin{solucion}
Si $D$ es denso y numerable, podemos tomar $L=D$, pues $D\subset\operatorname{span}D$ implica densidad del subespacio generado.

Para el recíproco, sea $\F_0=\Q$ en el caso real y $\F_0=\Q+i\Q$ en el complejo. Sea $D$ el conjunto de todas las combinaciones lineales finitas de elementos de $L$ con coeficientes en $\F_0$, incluyendo el vector cero. Para cada longitud de combinación hay solo un producto finito de conjuntos numerables; al tomar la unión sobre todas las longitudes, $D$ es numerable.

Probemos su densidad. Dado $v\in V$ y $\eps>0$, la densidad de $\operatorname{span}L$ permite escoger $w=\sum_{j=1}^m\alpha_jv_j$ con $\norm{v-w}<\eps/2$. Para cada $j$, elijamos $\beta_j\in\F_0$ tal que
\[
|\alpha_j-\beta_j|<\frac{\eps}{2m(1+\norm{v_j})}.
\]
Los racionales son densos en $\R$, y $\Q+i\Q$ es denso en $\C$ al aproximar por separado las partes real e imaginaria; por ello estas elecciones son posibles. Entonces $u=\sum_{j=1}^m\beta_jv_j\in D$ y
\[
\norm{w-u}\le\sum_{j=1}^m|\alpha_j-\beta_j|\norm{v_j}<\eps/2.
\]
La triangular da $\norm{v-u}<\eps$. Si la combinación es vacía, se toma $u=0$; en particular, $V=\{0\}$ queda cubierto. Así $D$ es denso.
\end{solucion}

\begin{ejercicio}[Lectura 6, ejercicio 3]\label{ej:lectura-6-3}
Si $(A_n)$ es una familia a lo sumo numerable de subespacios separables de un espacio métrico $X$, demostrar que $Y=\bigcup_nA_n$ es separable.
\end{ejercicio}
\begin{solucion}
Para cada $A_n$ no vacío, elijamos $D_n\subset A_n$ denso en su topología relativa y numerable; para el vacío tomemos $D_n=\varnothing$. El conjunto $D=\bigcup_nD_n$ es numerable. Si $y\in Y$ y $r>0$, existe $n$ con $y\in A_n$. Por densidad relativa, $B_X(y,r)\cap A_n$ contiene un punto de $D_n$, y por tanto $B_Y(y,r)$ contiene un punto de $D$. Esto demuestra $\cl D^{\,Y}=Y$. No se necesita que los $A_n$ sean abiertos ni cerrados.
\end{solucion}

\begin{ejercicio}[Lectura 6, ejercicio 4]\label{ej:lectura-6-4}
Demostrar que todo subespacio de un espacio métrico separable es separable.
\end{ejercicio}
\begin{solucion}
Sea $\{B_n\}$ una base numerable de $X$, cuya existencia asegura el teorema~\ref{ch06:separablebase}. Los conjuntos $Y\cap B_n$ forman una base de $Y$: si $U=Y\cap G$ es abierto relativo y $y\in U$, hay $B_n$ con $y\in B_n\subset G$, y entonces $y\in Y\cap B_n\subset U$. Esta base es numerable. Elegimos un punto de cada elemento no vacío de ella; por el lema~\ref{ch06:baseadenso}, esos puntos forman un denso numerable de $Y$.

No basta intersectar con $Y$ un denso numerable prefijado de $X$. Por ejemplo, $\Q$ es denso en $\R$, pero su intersección con $Y=\R\setminus\Q$ es vacía. La prueba con una base elige puntos dentro del subespacio y evita ese error.
\end{solucion}

\begin{ejercicio}[Lectura 6, ejercicio 5]\label{ej:lectura-6-5}
Sean $X_1,\ldots,X_m$ espacios métricos no vacíos. Demostrar que su producto finito, con una métrica producto, es separable si y solo si lo es cada factor.
\end{ejercicio}
\begin{solucion}
Podemos usar $d(x,y)=\max_{1\le j\le m}d_j(x_j,y_j)$. Su bola de radio $r$ es el producto de las bolas de radio $r$ de los factores. Si $D_j$ es denso numerable en $X_j$, el producto $D=\prod_{j=1}^mD_j$ es numerable. Dado $x$ y $r>0$, elegimos $a_j\in D_j$ con $d_j(x_j,a_j)<r$ para cada uno de los finitísimos índices. Entonces $a\in D\cap B(x,r)$ y $D$ es denso.

Recíprocamente, sea $D$ denso numerable en el producto. La proyección $\pi_j(D)$ es numerable. Dado $x_j\in X_j$ y $r>0$, elegimos puntos arbitrarios en los demás factores y formamos $x$ con coordenada $j$ igual a $x_j$. Una aproximación $a\in D$ con $d(a,x)<r$ proporciona $d_j(a_j,x_j)<r$, así que $\pi_j(D)$ es denso en $X_j$.

Las métricas producto finitas dadas por suma, máximo o norma $p$ inducen la misma topología: para números no negativos $t_j$, $\max_jt_j\le(\sum_jt_j^p)^{1/p}\le m^{1/p}\max_jt_j$ y $\max_jt_j\le\sum_jt_j\le m\max_jt_j$. Estas cotas comparan las bolas y trasladan el resultado a cualquiera de esas métricas. La hipótesis de factores no vacíos es necesaria: un factor vacío hace que el producto sea vacío, aunque otro factor sea no separable.
\end{solucion}

\begin{definicion}\index{denso en ninguna parte}
Un conjunto $A\subset X$ es denso en ninguna parte si $\operatorname{int}(\cl A)=\varnothing$. La clausura dentro de esta definición es esencial: tener interior vacío por sí solo no basta, como muestra $\Q\subset\R$.
\end{definicion}

\begin{ejercicio}[Lectura 6, ejercicio 6]\label{ej:lectura-6-6}
Demostrar que un conjunto denso en ninguna parte tiene complemento denso, y que un conjunto denso y cerrado tiene complemento denso en ninguna parte.
\end{ejercicio}
\begin{solucion}
Si $\operatorname{int}(\cl A)=\varnothing$, ningún abierto no vacío $U$ está contenido en $\cl A$. Por tanto $U\cap(X\setminus\cl A)\ne\varnothing$, y en particular $U\cap(X\setminus A)\ne\varnothing$. Así el complemento de $A$ es denso. Hemos probado algo más fuerte: el complemento de la clausura es abierto y denso.

Si $A$ es denso y cerrado, $A=\cl A=X$; por tanto su complemento es vacío, cuya clausura tiene interior vacío. La segunda afirmación de la lectura es correcta, aunque trivial con ambas hipótesis. Una formulación adicional útil es que el complemento de un abierto denso es cerrado y denso en ninguna parte: si $G$ es abierto y denso, $F=X\setminus G$ es cerrado y no contiene ningún abierto no vacío, porque todo abierto no vacío intersecta a $G$. Así $\operatorname{int}(\cl F)=\operatorname{int}F=\varnothing$.
\end{solucion}

\begin{ejercicio}[Lectura 6, ejercicio 7]\label{ej:lectura-6-7}
Demostrar que en un espacio métrico no vacío un conjunto no puede ser simultáneamente denso y denso en ninguna parte.
\end{ejercicio}
\begin{solucion}
Si $A$ es denso, $\cl A=X$. Entonces $\operatorname{int}(\cl A)=\operatorname{int}X=X\ne\varnothing$, lo que contradice ser denso en ninguna parte. La condición de espacio no vacío evita la excepción $X=A=\varnothing$.
\end{solucion}

\begin{ejercicio}[Lectura 6, ejercicio 8]\label{ej:lectura-6-8}\label{ch06:intervalos}
Demostrar que todo abierto de $\R$ es unión a lo sumo numerable de intervalos abiertos disjuntos dos a dos.
\end{ejercicio}
\begin{solucion}
Sea $U\subset\R$ abierto. Si es vacío, la unión vacía resuelve el problema. Para $x\in U$, definamos $I_x$ como la unión de todos los intervalos abiertos contenidos en $U$ que contienen a $x$. Hay al menos uno porque $U$ es abierto. La unión es abierta y está contenida en $U$.

Además, $I_x$ es un intervalo: si $a<b$ pertenecen a $I_x$, hay intervalos $J_a,J_b$ de la familia que contienen respectivamente a $a,b$ y ambos contienen a $x$. Como se intersectan, su unión es un intervalo. Para justificar este último hecho, un punto entre dos puntos de $J_a\cup J_b$ queda en uno de los intervalos, o entre ellos; en el segundo caso un punto común a ambos permite dividir el segmento en dos partes contenidas en los respectivos intervalos. Por ello todos los puntos entre $a,b$ pertenecen a $J_a\cup J_b\subset I_x$.

Así $I_x$ es un intervalo abierto, posiblemente una semirrecta o toda $\R$. Es maximal entre los intervalos de $U$ que contienen a $x$, porque su definición ya reúne a todos ellos. Si $I_x$ e $I_y$ se intersectan, su unión es un intervalo abierto contenido en $U$; contiene tanto a $x$ como a $y$. Por maximalidad está contenida tanto en $I_x$ como en $I_y$, y por tanto $I_x=I_y$. Los intervalos distintos de esta colección son disjuntos y cubren $U$ porque $x\in I_x$.

Cada intervalo no vacío contiene un racional. Fijemos una enumeración de $\Q$ y asignemos a cada intervalo el primer racional de la lista que contiene. Dos intervalos disjuntos reciben racionales distintos, de modo que la familia se inyecta en $\Q$ y es a lo sumo numerable. Esto concluye la descomposición, incluidos los intervalos no acotados.
\end{solucion}

\section{Una misma precisión para todos los puntos}\label{lectura:07}
La densidad explica qué puntos bastan para aproximar el espacio; ahora examinaremos qué control permite transportar esas aproximaciones a través de una función. La continuidad ordinaria permite elegir la precisión de entrada después de fijar un punto. La continuidad uniforme exige una elección que funcione simultáneamente en todo el dominio. Esta diferencia de cuantificadores explica por qué una función puede comportarse bien cerca de cada punto y, sin embargo, amplificar mucho las distancias en regiones cada vez más lejanas. Usaremos continuidad, sucesiones y la desigualdad triangular; no necesitaremos compacidad.

\begin{definicion}\label{ch07:def-uniforme}
Una aplicación $f:(X,d)\to(Y,\rho)$ es \emph{uniformemente continua} si
\[
 \forall\eps>0\ \exists\delta>0\ \forall x,y\in X,
 \qquad d(x,y)<\delta\Longrightarrow\rho(f(x),f(y))<\eps.
\]
Es \emph{Lipschitz} con constante $L\geq0$ si
\[
 \rho(f(x),f(y))\leq Ld(x,y)\qquad(x,y\in X).
\]
\end{definicion}
\index{continuidad!uniforme}\index{Lipschitz, funcion de@Lipschitz, función de}
La desigualdad de Lipschitz es no estricta: escribir $<$ para todo par de puntos produciría $0<0$ al tomar $x=y$. Una constante de Lipschitz no tiene por qué ser óptima; cualquier constante mayor también sirve.

\begin{proposicion}\label{ch07:lip-uniforme}
Toda aplicación Lipschitz es uniformemente continua, y toda aplicación uniformemente continua es continua.
\end{proposicion}
\begin{proof}
Si $L>0$, para una tolerancia $\eps>0$ tomamos $\delta=\eps/L$. Si $L=0$, todas las imágenes coinciden y sirve cualquier $\delta>0$. En ambos casos se cumple la definición. Para la segunda afirmación fijamos $x\in X$: el mismo $\delta$ que funciona para todos los pares funciona, en particular, para los pares $(x,y)$. Ésta es precisamente la continuidad en $x$.
\end{proof}

\begin{ejemplo}[Qué significa aquí discreto]
Si $d(x,y)=1$ para $x\ne y$, toda aplicación de $X$ a cualquier espacio métrico es uniformemente continua: con $\delta=1$ la condición $d(x,y)<\delta$ obliga a $x=y$. El argumento vale siempre que exista $c>0$ con $d(x,y)\geq c$ para puntos distintos; se dice entonces que $X$ es uniformemente discreto. Tener solamente topología discreta no basta. En $X=\{1/n:n\in\N\}\subset\R$, todos los puntos son aislados, pero $f(1/n)=n$ no es uniformemente continua: las distancias entre $1/n$ y $1/(n+1)$ tienden a cero, mientras que las imágenes distan $1$.
\end{ejemplo}

\begin{ejemplo}[Continuidad sin uniformidad]\label{ch07:cuadrado}
La función $f(x)=x^2$ es continua en $\R$, pero no uniformemente continua. Tomemos $x_n=n$ y $y_n=n+1/n$. Entonces
\[
 |x_n-y_n|=1/n\longrightarrow0,
 \qquad |f(x_n)-f(y_n)|=2+1/n^2\geq2.
\]
Si existiera un $\delta$ uniforme para $\eps=1$, un $n$ con $1/n<\delta$ contradiría la segunda desigualdad. Este razonamiento fija primero la tolerancia que fracasa y después produce pares arbitrariamente próximos. No es válido deducir una cota inferior proporcional a $\delta$ a partir de la sola desigualdad $|x-y|<\delta$.
\end{ejemplo}

\begin{proposicion}[Composición]\label{ch07:composicion}
Si $f:X\to Y$ y $g:Y\to Z$ son uniformemente continuas, también lo es $g\circ f$.
\end{proposicion}
\begin{proof}
Fijado $\eps>0$, la uniformidad de $g$ proporciona $\eta>0$ tal que $\rho(u,v)<\eta$ implica $\sigma(g(u),g(v))<\eps$. Aplicamos ahora la uniformidad de $f$ con tolerancia $\eta$: existe $\delta>0$ tal que $d(x,y)<\delta$ implica $\rho(f(x),f(y))<\eta$. Las dos implicaciones encadenadas dan el resultado. El orden de elección es $\eps$, después $\eta$ y finalmente $\delta$.
\end{proof}

\section{Sucesiones cuyos términos terminan estando próximos}
\begin{definicion}\label{ch07:cauchy}
Una sucesión $(x_n)$ de un espacio métrico es \emph{de Cauchy} si
\[
 \forall\eps>0\ \exists N\in\N\ \forall n,m\geq N,
 \qquad d(x_n,x_m)<\eps.
\]
\end{definicion}
\index{Cauchy, sucesion de@Cauchy, sucesión de}
La lectura original utiliza la condición $d(x_n,x_N)<\eps$ para $n\geq N$. Ambas formulaciones son equivalentes: la definición anterior implica esa condición al tomar $m=N$; a la inversa, aplicamos la condición con $\eps/2$ y usamos
\[
 d(x_n,x_m)\leq d(x_n,x_N)+d(x_N,x_m)<\eps.
\]
La condición de Cauchy compara términos de la sucesión entre sí; no presupone que su posible límite pertenezca al espacio.

\begin{teorema}\label{ch07:cauchy-imagen}
Las aplicaciones uniformemente continuas transforman sucesiones de Cauchy en sucesiones de Cauchy.
\end{teorema}
\begin{proof}
Sea $\eps>0$ y elijamos $\delta>0$ por la continuidad uniforme de $f$. Como $(x_n)$ es de Cauchy, existe $N$ tal que $n,m\geq N$ implica $d(x_n,x_m)<\delta$. Por tanto $\rho(f(x_n),f(x_m))<\eps$ para esos índices, que es la condición de Cauchy en $Y$.
\end{proof}
La continuidad ordinaria no basta: $x_n=1/n$ es de Cauchy en $(0,1]$, pero la función continua $x\mapsto1/x$ la transforma en $(n)$, que no es de Cauchy.

\section{Ejercicios sobre continuidad uniforme y extensiones}
\begin{ejercicio}[Lectura 7, ejercicio 1]\label{ej:lectura-7-1}
Demuestre que la composición de dos funciones Lipschitz es Lipschitz y determine una constante para la composición.
\end{ejercicio}
\begin{solucion}
Si $f:X\to Y$ tiene constante $L$ y $g:Y\to Z$ tiene constante $M$, entonces
\[
 \sigma(g(f(x)),g(f(y)))\leq M\rho(f(x),f(y))\leq MLd(x,y).
\]
Así, $ML$ es una constante válida, incluso cuando alguno de los factores es cero. No necesariamente es la mejor: si $g$ es constante en $f(X)$, la composición tiene constante óptima cero aunque $L,M$ sean positivos.
\end{solucion}

\begin{ejercicio}[Lectura 7, ejercicio 2]\label{ej:lectura-7-2}
Sean $f:X\to Y$ y $g:Y\to Z$ continuas. Determine si $g\circ f$ debe ser uniformemente continua cuando: \begin{enumerate}[label=(\roman*)]
\item $f$ es uniformemente continua;
\item $g$ es uniformemente continua.
\end{enumerate}
\end{ejercicio}
\begin{solucion}
En ambos casos la respuesta es negativa. En (i) tomamos $X=Y=Z=\R$, $f(x)=x$ y $g(x)=x^2$. La identidad es $1$-Lipschitz y la composición es el cuadrado, que falla por el ejemplo~\ref{ch07:cuadrado}. En (ii) intercambiamos los papeles: $f(x)=x^2$ y $g(x)=x$. Ahora $g$ es $1$-Lipschitz y la composición vuelve a ser $x^2$.
\end{solucion}

\begin{ejercicio}[Lectura 7, ejercicio 3]\label{ej:lectura-7-3}
Sea $A$ un cerrado no vacío de un espacio métrico $X$. Pruebe que $x\mapsto\dist(x,A)$ es uniformemente continua.
\end{ejercicio}
\begin{solucion}
Para $a\in A$, la desigualdad triangular da $d(x,a)\leq d(x,y)+d(y,a)$. Tomando ínfimos en $a$ obtenemos
\[
 \dist(x,A)\leq d(x,y)+\dist(y,A).
\]
Al intercambiar $x,y$ concluimos
\[
 |\dist(x,A)-\dist(y,A)|\leq d(x,y).
\]
La función es $1$-Lipschitz. La prueba no exige que $A$ sea cerrado; esa hipótesis solamente permite identificar $\dist(x,A)=0$ con $x\in A$. Se agrega la condición $A\ne\varnothing$ porque la distancia al vacío no es una función real finita.
\end{solucion}

\begin{ejercicio}[Lectura 7, ejercicio 4]\label{ej:lectura-7-4}
Dé una función uniformemente continua entre espacios métricos que no sea abierta.
\end{ejercicio}
\begin{solucion}
La función constante $f:\R\to\R$, $f(x)=0$, es $0$-Lipschitz. Sin embargo, la imagen del abierto $(-1,1)$ es $\{0\}$, que no es abierto en $\R$. La continuidad controla preimágenes de abiertos; ser abierta controla imágenes y es una condición distinta.
\end{solucion}

\begin{ejercicio}[Lectura 7, ejercicio 5]\label{ej:lectura-7-5}\label{ch07:ej-extension}
Sea $A$ abierto en un espacio métrico $X$ y sea $f:X\to\R$ uniformemente continua en $A$. Demuestre que $f|_A$ se extiende continuamente a $\cl A$. Dé un contraejemplo si se reemplaza $\R$ por un espacio métrico arbitrario.
\end{ejercicio}
\begin{solucion}
Probaremos la afirmación más fuerte en que $A$ es cualquier subconjunto y el espacio de llegada $Y$ es completo. Si $A=\varnothing$, la extensión tiene dominio vacío y no hay nada que definir. Para $x\in\cl A$ elegimos $a_n\in A$ con $d(a_n,x)<1/n$. Esta sucesión converge y es de Cauchy. Por el teorema~\ref{ch07:cauchy-imagen}, $(f(a_n))$ es de Cauchy; la completitud de $Y$ proporciona un límite. Definimos
\[
 g(x)=\lim_{n\to\infty}f(a_n).
\]
La elección no afecta al valor: si $b_n\to x$ en $A$, entonces $d(a_n,b_n)\leq d(a_n,x)+d(x,b_n)\to0$; la uniformidad implica $\rho(f(a_n),f(b_n))\to0$. Tomando límites, los dos valores propuestos coinciden. Si $x\in A$, podemos elegir la sucesión constante $a_n=x$, por lo que $g(x)=f(x)$.

Además, $g$ es uniformemente continua. Dado $\eps>0$, elijamos $\delta>0$ de modo que $d(a,b)<\delta$ implique $\rho(f(a),f(b))<\eps/2$ para $a,b\in A$. Si $x,y\in\cl A$ y $d(x,y)<\delta/3$, elegimos $a_n\to x$, $b_n\to y$. Para $n$ suficientemente grande, $d(a_n,b_n)<\delta$; al tomar límites queda $\rho(g(x),g(y))\leq\eps/2<\eps$. La extensión es única porque dos aplicaciones continuas que coinciden en un conjunto denso coinciden en su clausura: evaluamos ambas sobre una sucesión $a_n\to x$.

Para el contraejemplo tomamos $X=\R$, $A=(0,1)$ y $Y=(0,1)$ con la distancia usual. Definimos $f(x)=x$ en $A$ y $f(x)=1/2$ fuera de $A$. Su restricción a $A$ es $1$-Lipschitz. Una extensión continua $g:[0,1]\to(0,1)$ tendría $g(0)=\lim_{n\to\infty}g(1/(n+2))=0$, imposible en el espacio de llegada. Lo que falla es su completitud.
\end{solucion}

\begin{ejercicio}[Lectura 7, ejercicio 6]\label{ej:lectura-7-6}
Para subconjuntos no vacíos $A,B$ de $X$, defina $d(A,B)=\inf\{d(a,b):a\in A,b\in B\}$. Pruebe que $f:X\to Y$ es uniformemente continua si y sólo si
\[
 d(A,B)=0\Longrightarrow\rho(f(A),f(B))=0
\]
para todos esos conjuntos.
\end{ejercicio}
\begin{solucion}
Si $f$ es uniformemente continua, dado $\eps>0$ elegimos el correspondiente $\delta>0$. La igualdad $d(A,B)=0$ permite encontrar $a\in A,b\in B$ con $d(a,b)<\delta$. Entonces $\rho(f(A),f(B))\leq\rho(f(a),f(b))<\eps$. Un número no negativo menor que todo $\eps>0$ es cero.

Para el recíproco hay que cuidar las \emph{distancias cruzadas}: saber que $\rho(f(x_n),f(y_n))\geq\eps_0$ no basta por sí solo para separar los conjuntos de imágenes. Si $f$ no es uniformemente continua, existen $\eps_0>0$ y pares $x_n,y_n$ tales que
\[
 d(x_n,y_n)<1/n,\qquad \rho(p_n,q_n)\geq\eps_0,
 \quad p_n=f(x_n),\ q_n=f(y_n).
\]
Pongamos $r=\eps_0/8$. Si alguna bola $B(z,r)$ contiene infinitos $p_n$, retenemos esos índices $I$. Para $i,j\in I$, tenemos $\rho(q_j,z)\geq\eps_0-r$ y, por tanto, $\rho(p_i,q_j)>\eps_0-2r$. Si alguna bola contiene infinitos $q_n$, el argumento con los papeles intercambiados da la misma separación.

Queda el caso en que toda bola de radio $r$ contiene sólo un número finito de términos de cada sucesión, contando sus índices. Elegimos inductivamente índices crecientes $n_k$. Una vez elegidos los primeros $k$, descartamos todos los índices $n$ para los que $p_n\in B(q_{n_j},r)$ o $q_n\in B(p_{n_j},r)$ para algún $j\leq k$. Sólo descartamos una unión finita de conjuntos finitos; por ello queda un índice mayor que $n_k$. La sucesión retenida satisface $\rho(p_{n_i},q_{n_j})\geq r$ si $i\ne j$, y la misma desigualdad vale cuando $i=j$ porque $\eps_0\geq r$.

En cualquiera de los casos obtenemos un conjunto infinito de índices $I$ y una constante $c>0$ tales que $\rho(p_i,q_j)\geq c$ para todo $i,j\in I$. Para $A=\{x_i:i\in I\}$ y $B=\{y_i:i\in I\}$ se cumple $d(A,B)=0$, pues $I$ no está acotado y $d(x_i,y_i)<1/i$, pero $\rho(f(A),f(B))\geq c$. Esto contradice la propiedad supuesta.
\end{solucion}

\begin{ejercicio}[Lectura 7, ejercicio 7]\label{ej:lectura-7-7}
Suponga que $f:X\to Y$ satisface $\rho(f(x),f(y))\geq\alpha d(x,y)$ para alguna $\alpha>0$. Pruebe que es inyectiva, que su inversa sobre la imagen es Lipschitz y que $f$ es abierta si $f(X)$ es abierto en $Y$.
\end{ejercicio}
\begin{solucion}
Si $f(x)=f(y)$, la desigualdad da $0\geq\alpha d(x,y)$ y por tanto $x=y$. Para $u=f(x),v=f(y)$ obtenemos
\[
 d(f^{-1}(u),f^{-1}(v))\leq\alpha^{-1}\rho(u,v).
\]
Así, $f^{-1}:f(X)\to X$ es continua. Si $U\subset X$ es abierto, $f(U)=(f^{-1})^{-1}(U)$ es abierto en el subespacio $f(X)$. Cuando $f(X)$ es abierto en $Y$, todo abierto suyo es abierto en $Y$, y $f$ es una aplicación abierta. La hipótesis no afirma que $f$ sea continua; controla únicamente su inversa.
\end{solucion}

\section{Continuidad, aproximación y control uniforme}
El criterio de tolerancias permite obtener cotas concretas; el de abiertos simplifica composiciones y pegados; el secuencial permite detectar discontinuidades mediante aproximaciones explícitas. Un homeomorfismo conserva esta estructura de continuidad en ambas direcciones. No tiene por qué conservar distancias, acotación ni propiedades que dependan de la métrica concreta.

En un espacio métrico, una familia numerable de puntos de aproximación se convierte en una familia numerable de entornos básicos. Esa conversión permite reducir cubiertas abiertas a subfamilias numerables y heredar separabilidad a todos los subespacios. La condición «denso en ninguna parte» mide una limitación distinta: ni siquiera después de agregar puntos límite aparece una región abierta completa.

La continuidad uniforme controla simultáneamente todos los puntos, preserva la condición de Cauchy y permite extender funciones a puntos límite cuando el espacio de llegada contiene los límites necesarios. El capítulo~\ref{chap:08} estudia precisamente esa propiedad del espacio de llegada: la completitud.
